\section{Theoretische Grundlagen}
\subsection{Smart Grids als komplexe Netzwerke}
Wie in der Einleitung bereits geschrieben stellen Smart Grids eine Weiterentwicklung konventioneller Stromnetze dar. Sie können als cyber-physische-Systeme , in denen physische Komponenten wie Kraftwerke, Umspannwerke oder Leitungen eng mit digitalen Kommunikations- und Steuerungseinheiten interagieren verstanden werden. Entscheidungen über Lastverteilung, Einspeisung oder Netzstabilisierung erfolgen zunehmend automatisiert und datengetrieben, wodurch neue Abhängigkeiten zwischen räumlich und funktional getrennten Netzkomponenten entstehen. Störungen können sich daher nicht nur entlang physischer Leitungen ausbreiten, sondern auch über informationstechnische Kanäle verstärkt oder ausgelöst werden.
Zur Analyse dieser komplexen Wechselwirkungen bietet sich eine Abstraktion von Smart Grids als Netzwerke an. In dieser Darstellung werden physische oder logische Komponenten des Stromnetzes als Knoten modelliert, während elektrische Leitungen, Kommunikationsverbindungen oder funktionale Abhängigkeiten als Kanten beschrieben werden. Eine solche graphbasierte Modellierung erlaubt es, strukturelle Eigenschaften des Gesamtsystems systematisch zu erfassen und vergleichbar zu machen. Insbesondere ermöglicht sie die Untersuchung, wie die Anordnung und Vernetzung der Komponenten die Stabilität, Fehlertoleranz und Resilienz des Netzes beeinflussen.
Die Anwendung der Theorie komplexer Netzwerke ist dabei nicht als Vereinfachung realer Stromnetze zu verstehen, sondern als analytischer Rahmen, der gezielt jene Aspekte hervorhebt, die für das Verhalten vernetzter Systeme unter Störungen entscheidend sind. Resilienz in Smart Grids wird maßgeblich durch die Struktur der Verbindungen bestimmt, über die sich Ausfälle, Überlastungen oder Fehlfunktionen ausbreiten oder abgefedert werden können. Netzwerkbasierte Ansätze ermöglichen es somit, strukturelle Ursachen von Verwundbarkeit zu identifizieren, ohne sich ausschließlich auf detaillierte physikalische Simulationsmodelle stützen zu müssen.
Gerade im Kontext zunehmender Dezentralität und Komplexität bieten netzwerktheoretische Methoden einen wichtigen Mehrwert, da sie den Fokus auf systemische Eigenschaften legen. Sie erlauben es, kritische Komponenten zu identifizieren, strukturelle Schwachstellen aufzudecken und allgemeine Aussagen über die Robustheit des Netzwerks gegenüber unterschiedlichen Störungsszenarien zu treffen. Damit bilden sie eine zentrale theoretische Grundlage für die in dieser Arbeit verfolgte Resilienzanalyse von Smart Grids.

\subsection{Netzwerkmodelle und strukturelle Eigenschaften}
Die Struktur eines Netzwerks hat einen entscheidenden Einfluss darauf, wie robust oder verwundbar es gegenüber Störungen ist. Um diese Zusammenhänge systematisch zu untersuchen, wurden in der Netzwerktheorie verschiedene Modellklassen entwickelt, die charakteristische Topologien realer Systeme abstrahieren. Ergänzend dazu ermöglichen strukturelle Netzwerkmetriken eine quantitative Beschreibung dieser Topologien und dienen als Erklärgrößen für das Verhalten von Netzwerken unter Ausfällen oder Angriffen.
\subsubsection{Netzwerktopologien}
Komplexe Netzwerke können auf unterschiedliche Weise strukturiert sein. Drei grundlegende Modelle haben sich bisher zur Beschreibung und Simulation realer Systeme etabliert: Erdős-Rényi-Netzwerke, Watts-Strogatz-Small-World-Netzwerke und Barabási-Albert-skalenfreie Netzwerke. Diese Topologien unterscheiden sich hinsichtlich Verbindungsstruktur, Gradverteilung, Clusterkoeffizient und typischer Pfadlängen und damit auch hinsichtlich ihrer Resilienz.
Das Erdős-Rényi-Netzwerk, ER-Modell, \citep{Erdos1959RandomGraphs} erzeugt ein Netzwerk, indem zwischen allen möglichen Knotenpaaren jeweils mit einer festen Wahrscheinlichkeit p eine Kante gezogen wird. Diese zufällige Verbindungsstrategie hat mehrere Konsequenzen:
	Die Gradverteilung ist binomial- bzw. in großen Netzen näherungsweise Poisson-verteilt und die meisten Knoten haben ähnliche Gradwerte.
	Der Clusterkoeffizient ist vergleichsweise gering, da nur wenige Dreiecksbeziehungen entstehen.
	Die Pfadlänge ist typischerweise klein (logarithmisch mit Netzwerkgröße).
	Die Resilienz ist moderat robust gegenüber zufälligen Ausfällen, aber vulnerabel gegenüber gezielten Angriffen, da keine systematische Struktur existiert, die kritische Knoten schützt.
ER-Netzwerke bilden eine wichtige Baseline, aber reale Stromnetze weisen häufig strukturierte Muster auf, die durch dieses Modell nicht vollständig beschrieben werden können.
Das Small-World-Netzwerk-Modell \citep{Watts1998SmallWorld} kombiniert Merkmale regulärer und zufälliger Netzwerke. Die Konstruktion erfolgt durch Ausgang eines regulären Rings, in dem Knoten mit ihren nächsten Nachbarn verbunden sind, und anschließend eine zufällige „Umverteilung" einzelner Kanten passiert.
Charakteristische Eigenschaften:
	Eine hohe lokale Clusterbildung - also eine lokale Redundanz - wird geschaffen. Diese Knoten haben viele gemeinsame Nachbarn.
	Die durchschnittliche Pfadlänge ist eher kurz. Dadurch existiert eine effiziente Erreichbarkeit im gesamten Netzwerk.
	Die Resilienz ist durch hohe lokale Clusterung vergleichsweise hoch und damit das Netzwerk widerstandsfähiger gegenüber Ausfällen. Jedoch ist es anfällig für Störungen, wenn die wenigen „Shortcut"-Kanten ausfallen, die die Pfadlänge stark reduzieren.
Small-World-Strukturen treten in vielen Energie- und Informationsnetzen auf, da sie Effizienz mit lokaler Stabilität verbinden.
Von besonderer Bedeutung sind skalenfreie Netzwerke nach Barabási und Albert, deren Knotengradverteilung einem Potenzgesetz folgt. Das BA-Modell \citep{Albert2000ErrorAttack} basiert auf „präferentieller Anbindung" von Knoten. Es werden nacheinander Knoten hinzugefügt und sie werden bevorzugt mit bereits gut vernetzten Knoten verbunden („the rich get richer"). Dadurch entstehen sog. Hubs.
Eigenschaften:
	Die Gradverteilung folgt einem Potenzgesetz (scale-free). Dadurch entstehen wenige Hubs mit vielen Knoten mit geringem Grad.
	Der Clusterkoeffizient ist vergleichsweise gering, kann aber je nach Erweiterung variieren.
	Die Pfadlängen ist sehr kurz, da Hubs viele Wege verbinden (ultra-small world).
	Die Resilienz ist hoch, da das Netzwerk sehr robust gegen zufällige Ausfälle ist (da diese meist kleine Knoten betreffen). Es ist aber stark anfällig gegenüber gezielten Angriffen auf Hubs, da das Netzwerk leicht fragmentieren kann.
Viele reale Infrastrukturnetze, einschließlich Kommunikationssysteme, Internet-Backbones und Teile moderner Stromnetze, zeigen skalenfreie Eigenschaften. Daher ist das BA-Modell besonders relevant für Analysen der Verwundbarkeit kritischer Infrastrukturen.

\subsubsection{Strukturelle Netzwerkmetriken}

Um die Robustheit und Verwundbarkeit eines Netzwerks quantitativ zu beschreiben, werden verschiedene strukturelle Metriken herangezogen. Eine zentrale Rolle spielt dabei die Gradverteilung.
Die Gradverteilung $P(k)$ ist die Wahrscheinlichkeit, dass ein Knoten eines Netzwerks Grad k hat und liegt zwischen 0 und 1. Die Gradverteilung eines Netzwerks zeigt die Beschaffenheit deren Gesamtstruktur auf und kann somit der Feststellung damit einhergehender, typischer Eigenschaften dienen. Für Smart Grids ist die Feststellung solcher typischen Eigenschaften insofern relevant, als dass hierdurch kritische Hubs identifiziert werden können, wie z.B. große Umspannwerke mit einer besonders hohen Anzahl an Leitungsverbindungen. Die Gradverteilung von Stromnetzwerken weist die grundsätzlichen Bedingungen der Skalenfreiheit auf: Zum Einen ist durch stetigen Ausbau des Stromnetzes das Wachstum als Voraussetzung gegeben. Ebenso werden die Kraftwerke vorzugsweise über stärker vernetzte Umspannwerke angebunden, anstatt peripher dezentral geleitet zu werden. Somit entsteht die typische Gradverteilung, die dem Potenzgesetz entspricht $P(k) ~ k-\gamma$. Bei dieser Verteilung wird die Wahrscheinlichkeit $P(k)$ bei wachsendem k immer kleiner. Im Falle von Stromnetzwerken, bspw. in der westlichen USA, liegt $\gamma$ ca. bei 4 und somit an der Grenze der Eigenschaften der Skalenfreiheit : die Anzahl großer Hubs ist vergleichsweise gering und stößt mit einer Gesamtzahl von einigen tausend Kraftwerken relativ schnell an eine natürliche Obergrenze. 
Der Clustering-Koeffizient C beschreibt, ob zwei Nachbarn eines Knotens ebenfalls füreinander Nachbarn sind und damit die Ausprägung der Transitivität des Netzwerkes : 
\[
C = \frac{\text{Anzahl geschlossener Pfade der L\"ange 2}}{\text{Anzahl Pfade der L\"ange 2}}
\]
Bei einer sehr großen Anzahl an Knoten N eines Netzwerks wird der Clustering-Koeffizient sehr klein. Diese Ausprägung eines Extremfalls ist bei Stromnetzen aufgrund der begrenzten Knotenanzahl jedoch nicht gegeben. Relevant wird der Clustering-Koeffizient dann, wenn einzelne Knoten und Verbindungen ausfallen. Der Koeffizient gibt an, wie groß die Redundanz im Netzwerk ist. Ist ein hoher Clustering-Koeffizient gegeben, so stehen im Netzwerk mehrere alternative Pfade bei Ausfällen zur Verfügung, was für Robustheit des Netzwerks gegen lokale Ausfälle sorgt. 
Die mittlere Pfadlänge beschreibt die mittlere kürzeste Entfernung zwischen Knoten. Dieser Kennwert ist relevant, um die Ausbreitung von Lastfluss oder Störungen zu beschreiben und um die Wiederherstellungszeit einzuschätzen . Die mittlere Pfadlänge lässt sich wie folgt berechen:
\[
\ell = \frac{1}{N(N-1)} \sum_{i j} d_{i j}
\]
Wobei $d_{ij}$ die durchschnittliche kürzeste Distanz zwischen den Knoten i und j eines Netzwerks bezeichnet.

\subsection{Robustheit, Resilienz, Fehlertoleranz}
In der Literatur zu komplexen Netzwerken werden die Begriffe Robustheit, Fehlertoleranz und Resilienz häufig gemeinsam verwendet, obwohl sie sich auf unterschiedliche Aspekte der Systemstabilität beziehen. In den folgenden Unterkapiteln wird eine Abgrenzung im Kontext von Stromnetzwerken vorgenommen, basierend auf der Übersichtsarbeit von \cite{Qi2024}. 
\subsubsection{Robustheit}
Robustheit kann als Teilaspekt von Resilienz betrachtet werden  und bezieht sich hauptsächlich darauf, wie widerstandsfähig ein Netzwerk in seiner Struktur selbst gegenüber Störungen ist. Im Bereich von Stromnetzwerken, die zu den technologischen Netzwerken zählen, wird Robustheit als die Resistenz gegenüber Ausfällen bezeichnet. Die Kernfrage der Robustheit besteht darin, wie gut die Netzwerkfunktion unmittelbar nach einer Störung erhalten bleibt. In diesem Sinne gilt ein Stromnetz als robust, der Ausfall von einzelnen Komponenten wie Umspannwerken oder Leitungen nicht zu einer großflächigen Fragmentierung führt und ein Großteil der Knoten weiterhin über zusammenhängende Pfade erreicht werden kann. 
Gemessen wird die Robustheit typischer Weise über die größte zusammenhängende Komponente (GCC). Hierbei werden nach und nach Knoten oder Kanten entweder zufällig oder gezielt aus dem Netzwerk entfernt und nach jedem Schritt die Größe des größten funktionierenden Netzteils gemessen: 
\[
\mathrm{R} = \frac{1}{N} \sum_{q=1}^{N} S(q)
\]
Wobei R die Gesamtrobustheit eines Netzwerks bezeichnet und S(q) die relative Größe der größten Komponente nach q entfernten Knoten darstellt. 

\subsubsection{Fehlertoleranz und Resilienz}
Die Fehlertoleranz bezeichnet die Fähigkeit eines Systems, trotz Fehler einzelner Komponenten weiterhin funktionsfähig zu bleiben, ohne die Fehler zwingend reparieren zu müssen. Während Stromnetze im Bereich der technologischen Netze angesiedelt sind, bezieht sich die Fehlertoleranz mehr auf Informationsnetzwerke, zu denen Kommunikations- und Internetnetzwerke zählen. Dabei wird stärker auf Hard- und Software Bezug genommen und die Fähigkeit beschrieben, Fehler strukturell oder algorithmisch zu kompensieren, z.B. durch Redundanzen oder Umleitung von Flüssen. Die Kernfrage der Fehlertoleranz ist, wie gut das System während eines auftretenden Fehlers weiterarbeiten kann. So gilt ein Stromnetz als fehlertolerant, wenn bei einem Leitungsausfall automatisch alternative Wege für den Stromfluss gewählt werden. 
Da für die Analyse der Fehlertoleranz eine Modellierung von Flussdynamiken und Steuerungslogiken erforderlich ist, wird diese in der vorliegenden Arbeit nicht explizit untersucht. Sie stellt jedoch einen wichtigen ergänzenden Aspekt zur strukturellen Robustheit in Smart Grids dar. 

Die Resilienz umfasst die Widerstandsfähigkeit eines Netzwerks im Störungsfall, die Fähigkeit zur Begrenzung der Auswirkungen und die Fähigkeit zur Wiederaufnahme eines akzeptablen Funktionsniveaus. Dementsprechend ist die Resilienz ein übergreifendes Konzept, welches die Robustheit, Fehlertoleranz, Anpassungs- und Erholungsfähigkeit beinhaltet. Ein Smart Grid ist beispielsweise dann resilient, wenn es nach einem extremen Wetterereignis schnell wieder auf ein normales Versorgungsniveau zurückkehren kann. 
Die Kernfrage der Resilienz ist, wie viel der ursprünglichen Leistung noch vorhanden ist. Die Bestimmung der Resilienz erfordert neben strukturellen Modellen auch dynamische Modelle zur Erholung, Reparatur und zur Systemsteuerung. Diese umfassende Betrachtung geht über den Fokus der vorliegenden Arbeit hinaus. 
Zusammenfassend lässt sich sagen, dass Robustheit, Fehlertoleranz und Resilienz verschiedene Perspektiven auf die Stabilität von Stromnetzen und Smart Grid Netzen bieten. Der Fokus dieser Arbeit liegt jedoch auf struktureller Robustheit, um zu analysieren, wie die Netzwerkstruktur sowie die Interdependenzen zwischen Kommunikations- und Stromnetz die Anfälligkeit gegenüber Angriffen beeinflussen.